\documentclass[12pt]{article}
\usepackage[utf8]{inputenc}
\usepackage[T1]{fontenc}
\usepackage[serbian]{babel}
\usepackage{indentfirst} 
\usepackage{geometry}
\geometry{a4paper, margin=2.5cm}

\begin{document}

\begin{center}
    {\huge \textbf{Ko je napisao Mesečevu sonatu -- srce Betovena ili glas umetničkog doba?}} \\[0.4cm]
    {\small \textit{Autor: Dorotea Borenović, diplomirani profesionalni plesač, BŠNS}} \\
    {\small \textit{Predmet: Istorija umetnosti}}
\end{center}

\vspace{1cm}

Ludvig van Betoven je jedan od najpoznatijih kompozitora u istoriji muzike, ali i jedan od onih čiji je život bio obeležen velikom ličnom tragedijom – postepenim gubitkom sluha. Za jednog muzičara, gluvoća je verovatno nešto nezamislivo, ali upravo u tom periodu Betoven stvara neka od svojih najdubljih i najemotivnijih dela. Mesečeva sonata se često doživljava kao muzička ispovest, delo puno tišine, nemira i unutrašnje borbe, što prirodno pokreće pitanje njene inspiracije.

Betoven stvara na prelazu iz klasicizma u romantizam. Iako je potekao iz klasične muzičke tradicije, kroz svoju muziku sve više pokazuje osobine romantičara: jake emocije, lični doživljaj, unutrašnje konflikte i udaljavanje od strogih pravila forme. Zbog toga su njegova rana dela često nailazila na kritike, jer su se razlikovala od onoga na šta je tadašnja publika bila naviknuta i što se smatralo „lepom“ i pravilnom muzikom. Generalno svaki umetnik - muzičar, književnik, slikar - je u periodu klasicizma morao da poštuje određena pravila. Bilo je gotovo nemoguće da neko odluči da svoju ideju, koja se možda razlikuje od normi klasicizma, predstavi javno baš zbog tih strogih pravila. Zato svaka slika iz klasicizma po malo liči jedna na drugu. Ideje su morale bite iste i pravila su se morala poštovati. Zato svako književno delo iz klasicizma ima istu formu - radnja u 24 sata, glavni lik mora da ostane u centru pažnje, mesto radnje ne može da se menja. U muzici se lako može primetiti da su kompozicije dosta slične i nije teško zaključiti da su pisana po strogo propisanim pravilima. Iz tog razloga Betoven se izuzetno ne uklapa u ovu grupu umetnika.

Moja ideja je – da li je Betoven, u periodu kada je postepeno gubio sluh, inspiraciju pronalazio u drugim umetnostima – pre svega u slikarstvu i književnosti. Da li je posmatrao umetnička dela ili čitao književne tekstove koji su u njemu budili snažne emocije, pa je te emocije pokušavao da prevede u muziku? U tom smislu, njegova gluvoća možda nije značila potpunu izolaciju od umetnosti, već drugačiji način doživljavanja sveta – kroz slike, reči, misli i unutrašnja osećanja.

Ovo pitanje mi se javlja iz jednostavne, ali veoma bitne dileme: kako neko ko ne čuje u potpunosti može da zna koliko je njegovo delo zaista „dobro“? Poznato je da Betoven u kasnijem periodu gotovo uopšte nije mogao da čuje zvukove, ali je ipak nastavio da komponuje. Postoje istorijski podaci da je koristio vibracije klavira – oslanjajući se na rezonancu i frekvencije koje je mogao da oseti preko tela, naročito tako što je dodirivao instrument ili čak stavljao drvenu letvu između svojih zuba i klavira kako bi uspeo da oseti vibracije zvuka. Ipak, takav način rada može objasniti tehničku kontrolu nad muzikom, ali teško može u potpunosti objasniti emotivnu dubinu njegovih dela - pogotovo Mesečve sonate koja je meni i probudila ovo pitanje.

Mesečeva sonata, napisana 1801. godine, nastaje u periodu kada Betoven još nije bio potpuno gluv, ali je već bio svestan da mu se sluh nepovratno pogoršava. Upravo zato se ovo delo često tumači kao izraz unutrašnje napetosti, tuge, borbe sa samim sobom i povučenosti. Njena sporost, tišina između tonova i gotovo meditativni karakter ukazuju na muziku koja ne želi da impresionira spolja, već da izrazi unutrašnje stanje kompozitora. Zbog toga se nameće pitanje da li je Betoven u ovom periodu stvarao više oslanjajući se na unutrašnje slike i emocije nego na sam zvuk.

Istovremeno, Mesečeva sonata se često povezuje sa Betovenovom ličnom emotivnom pričom i mogućom zaljubljenošću, što dodatno komplikuje razumevanje njene inspiracije. Da li je ovo delo prvenstveno rezultat lične patnje i neuzvraćenih osećanja, ili je nastalo pod uticajem umetničkog duha vremena u kojem se Betoven našao? Poznato je da je bio u kontaktu sa književnicima i intelektualcima svog doba i da je pokazivao veliko interesovanje za ideje romantizma, koji je naglašavao emociju, subjektivni doživljaj i unutrašnji svet pojedinca.

Zbog svega toga, čini se opravdanim postaviti pitanje da li je Betoven, suočen sa gubitkom sluha, inspiraciju sve više tražio van same muzike – u slikama, tekstovima i idejama drugih umetnika – i da li je upravo taj spoj unutrašnjih osećanja i umetničkih uticaja oblikovao jedno od njegovih najpoznatijih dela.

\paragraph{Mogući umetnički susreti i zajednički duh vremena}


\end{document}
